\begin{table}[tp]
\centering
\caption{Second-order nonconvex--PL minimax optimization: \href{https://github.com/k3nfalt/longtails/blob/main/LongTails/Setups/NonconvexPLSecondOrder.lean}{NonconvexPLSecondOrder}.}\label{tab:NonconvexPLSecondOrder}
\parbox{\linewidth}{\footnotesize On finite-dimensional real inner-product spaces, $f$ has an $L$-Lipschitz gradient and differentiable gradient with $\rho$-Lipschitz Hessian. Inner maxima are attained and $2\mu(\Phi(x)-f(x,y))\le\|\nabla_y f(x,y)\|^2$ uniformly. Assume $0<\mu\le L$, $\rho,\varepsilon>0$, $\Phi(x_0)-\inf\Phi\le\Delta$ and $\Phi(x_0)-f(x_0,y_0)\le D$, with $\Delta,D\ge0$. Exact queries return $(f,\nabla f,\nabla^2f)$. The target is $\|\nabla\Phi(x)\|\le\varepsilon$ and $\nabla^2\Phi(x)\succeq-\sqrt{100\rho\sqrt\kappa(4+\kappa)\varepsilon}\,I$; the constructed output is deterministic.}\par\smallskip
\begin{threeparttable}
\footnotesize
\setlength{\tabcolsep}{4pt}\renewcommand{\arraystretch}{1.25}
\begin{tabularx}{\linewidth}{@{}>{\raggedright\arraybackslash}p{0.20\linewidth}>{\raggedright\arraybackslash}p{0.32\linewidth}>{\raggedright\arraybackslash}p{0.25\linewidth}>{\raggedright\arraybackslash}X@{}}
\toprule
\textbf{Result} & \textbf{Oracle complexity} & \textbf{Novelty} & \textbf{Limitations} \\
\midrule
\textbf{(1)} \href{https://github.com/k3nfalt/longtails/blob/main/LongTails/Theorems/EnvelopeCubicNonconvexPL.lean}{Envelope\allowbreak Cubic\allowbreak Nonconvex\allowbreak PL} \textbf{(Ours)} $(\star\star\star)$ \href{https://github.com/k3nfalt/longtails/blob/main/scripts/reviews/EnvelopeCubicNonconvexPL.md}{(review)}
& $\mathcal{O}\!\left(\left[1+\kappa+\dfrac{\Delta\sqrt\rho\kappa^{3/4}}{\varepsilon^{3/2}}+\dfrac{\Delta\rho^2}{\mu^3}\right]\ell^2\right)$ joint calls. & & \\
\midrule[1.2pt]
\textbf{(2)} ECRN \citep{ncplZhou2026}, Thm.~11. Strongly concave inner subclass.
& $\mathcal{O}\!\left(\sqrt\kappa\,\ell+\left[1+\dfrac{\Delta\sqrt\rho\kappa^{3/4}}{\varepsilon^{3/2}}\right](1+\log\ell)\right)$ gradient and Hessian calls; $\varepsilon\le L^2/[\rho(4+\kappa)^{3/2}]$.
& \textbf{(1)}: Scope: PL permits nonconcave inner objectives and nonunique maximizers. Matches the leading polynomial rate on the strongly concave overlap.
& \textbf{(1)}: $\kappa$ initialization instead of $\sqrt\kappa$; larger logarithmic overhead, function values required, curvature constant up to $10\times$ looser. Retuning \textbf{(2)} removes its accuracy restriction. \\
\midrule
\textbf{(3)} MCN \citep{ncplLuo2022}, Cor.~1. Strongly concave inner subclass.
& $\mathcal{O}\!\left(\left[\sqrt\kappa+\dfrac{\Delta\sqrt\rho\kappa^2}{\varepsilon^{3/2}}\right]\ell\right)$ total gradient/Hessian calls; Hessians: $\mathcal{O}(1+\Delta\sqrt\rho\kappa^{3/2}\varepsilon^{-3/2})$.
& \textbf{(1)}: Scope: strong concavity to PL. On the overlap, for $\varepsilon\le L^2/(\rho\kappa^{3/2})$, leading cost replaces $\kappa^2$ by $\kappa^{3/4}$, up to logarithms; curvature scaling improves by $\kappa^{3/4}$.
& \textbf{(1)} requires values and has larger initialization/logarithmic overhead. Leading-term improvement need not dominate total cost. \\
\midrule
\textbf{(4)} Smoothed-AGDA \citep{ncplYang2022}, Thm.~4.1, Cor.~4.1. Full inner-PL class.
& $\mathcal{O}\!\left(1+\dfrac{\kappa L(\Delta+b_0)}{\varepsilon^2}+\kappa(1+\log\kappa)\right)$ gradient calls, including translation to $\|\nabla\Phi\|\le\varepsilon$.
& \textbf{(1)} adds second-order stationarity. Costs are not directly comparable: different oracle access and accuracy guarantees.
& \textbf{(1)} requires a Lipschitz Hessian and full-Hessian queries; \textbf{(4)} requires only gradient smoothness. \\
\bottomrule
\end{tabularx}
\begin{tablenotes}[flushleft]\footnotesize
\item[] \textbf{Abbr:} $\Phi(x)=\max_y f(x,y)$; $\kappa=L/\mu$; $\ell=\log(2+\kappa+\rho^2(\Delta+D)/\mu^3+\rho\varepsilon/\mu^2+\mu^2/(\rho\varepsilon))$. PL: Polyak--\L{}ojasiewicz inequality. $b_0=2[\max_y\widehat f(x_0,y)-\min_x\widehat f(x,y_0)]$, where $\widehat f(x,y)=f(x,y)+L\|x-x_0\|^2$, is the initial auxiliary primal--dual gap. Curvature tolerances in \textbf{(2)} and \textbf{(3)} are respectively $\sqrt{\rho(4+\kappa)^{3/2}\varepsilon}$ and $\kappa^{3/2}\sqrt{\rho\varepsilon}$.
\end{tablenotes}
\end{threeparttable}
\end{table}
\clearpage
